%% ============================================================
%% SECTION 7: CONCLUSION
%% ============================================================
\section{Conclusion}
\label{sec:conclusion}

\subsection{Summary of Contributions}
\label{sec:conc_summary}

This paper presented VABQ (Variance-Aware Adaptive Bit-Width Quantisation), a unified framework for energy-efficient indoor air quality monitoring that jointly optimises sensor-level ADC resolution and inference-level neural network precision. The framework addresses a critical gap in environmental IoT research: existing methods optimise sensor compression and model quantisation independently, leaving compound energy savings unrealised. By exploiting the temporal structure inherent in indoor pollutant signals---long stable periods interspersed with activity-driven transients---VABQ achieves energy reductions exceeding the sum of independent optimisation at each layer.

Four principal contributions were established. First, a dual-layer optimisation pipeline was developed that unifies per-channel adaptive ADC bit-width control (4, 8, or 16 bits) with inference-level INT8 post-training quantisation under a shared variance-entropy decision criterion. Second, a composite variance-entropy metric was introduced that combines signal magnitude variation with information-theoretic content, enabling robust discrimination between genuine transients and noise-driven fluctuations across diverse pollutant types. Third, comprehensive experimental validation on the DALTON dataset---spanning 30 deployment sites, 89.1 million sensor samples, and six months of continuous monitoring---demonstrated 61.3\% sensor-level energy reduction, 74.8\% inference-level energy reduction, and 96.1\% classification accuracy with only 1.2 percentage-point degradation relative to full-precision baseline. Fourth, systematic comparison against four baseline configurations and eight published methods confirmed that joint optimisation yields compound savings: the synergistic interaction between sensor bit reduction and PTQ effectiveness arises because narrower input dynamic ranges improve quantisation efficiency.

The experimental findings validate three key hypotheses. Temporal structure in indoor air quality signals is sufficient to enable adaptive bit-width allocation, with 52.3\% of samples allocated to 4-bit precision during stable periods. Joint sensor-inference optimisation yields synergistic energy savings, with combined accuracy loss (1.2 pp) less than the sum of independent losses (1.3 pp). The variance-entropy composite criterion provides robust control across diverse deployment contexts, maintaining accuracy above 95.5\% across a wide hyperparameter range ($\alpha \in [0.55, 0.75]$, $\theta_L \in [0.10, 0.20]$, $\theta_H \in [0.45, 0.65]$).

\subsection{Practical Impact and Applicability}
\label{sec:conc_impact}

For battery-powered IoT deployments employing duty-cycled operation, VABQ enables projected system-level energy savings of 51.6\% relative to uniform 16-bit sensing with FP32 inference. This translates to extending operational lifetime from approximately 72 hours to 148 hours on a standard 5000 mAh battery, enabling multi-day autonomous monitoring without recharging. For mains-powered building-scale deployments with hundreds of sensor nodes, the aggregate electricity cost and carbon footprint reductions are proportionally significant. At commercial electricity rates (\$0.15 per kWh), a 100-node deployment operating continuously for one year would save approximately \$47 annually per building through VABQ optimisation, with corresponding CO$_2$ emission reductions of 28 kg per building assuming grid carbon intensity of 0.5 kg CO$_2$ per kWh.

The framework's lightweight design---45,000-parameter 1D-CNN, $O(1)$ per-sample overhead, zero inter-node communication---ensures compatibility with resource-constrained edge hardware. The ESP32 microcontroller (520 KB SRAM, 240 MHz dual-core) executes VABQ with 5\% computational overhead, demonstrating feasibility for deployment on commodity IoT platforms. Calibration requires only one week of on-site data per deployment, substantially less burdensome than quantisation-aware training or supervised fine-tuning. These characteristics position VABQ as a practical, deployable solution for large-scale environmental monitoring networks spanning residential, commercial, and institutional settings.

Beyond air quality monitoring, the variance-entropy adaptive quantisation principle generalises to adjacent domains exhibiting temporal structure: water quality monitoring (dissolved oxygen, pH, turbidity), soil moisture sensing for precision agriculture, industrial condition monitoring (vibration, temperature, acoustic emission), and building energy management (HVAC optimisation, occupancy detection). The core insight---that exploiting signal stationarity at both sensing and inference layers yields compound savings---applies broadly to time-series IoT applications where energy efficiency is paramount.

\subsection{Research Outlook and Open Questions}
\label{sec:conc_outlook}

Three high-priority research directions emerge from this work. First, \textbf{hardware-in-the-loop experimental validation} is essential to transition from simulation-based energy estimation to measured power consumption. Implementing VABQ on ESP32 hardware with cycle-accurate current measurement (e.g., via INA226 precision shunt) would validate the 61.3\% ADC energy reduction, quantify reconfiguration overhead, and assess robustness to ADC non-linearities at low bit-widths. Such validation would also enable investigation of advanced power management strategies, including per-channel ADC clock gating and dynamic voltage-frequency scaling coordinated with bit-width allocation.

Second, \textbf{cross-dataset and cross-domain generalisation} must be established through systematic evaluation across diverse environmental monitoring applications. Immediate priorities include: (1) validation on the COLLECTiEF European indoor air quality dataset to assess transferability to buildings with advanced HVAC control; (2) extension to outdoor air quality networks (e.g., PurpleAir, OpenAQ) incorporating meteorological features and traffic patterns; (3) adaptation to water quality monitoring with chemical sensors (dissolved oxygen, conductivity, turbidity) exhibiting different drift characteristics; (4) deployment in industrial settings (manufacturing, mining, petrochemical) with extreme environmental conditions (temperature, humidity, vibration, electromagnetic interference). Cross-domain validation would identify necessary modifications to the variance-entropy criterion and quantify the range of applicability beyond indoor residential environments.

Third, \textbf{online adaptation and continual learning} are critical for long-term deployment (years) under sensor drift and evolving occupancy patterns. Current VABQ assumes stationary normalisation constants ($\sigma_{c,\max}^2$, $H_{c,\max}$) calibrated during commissioning. However, low-cost environmental sensors exhibit calibration drift (1--5\% per month)~\cite{dimitriou2025innovations}, and building occupancy patterns evolve (seasonal schedules, renovations, tenant turnover). Incremental learning methods that update normalisation constants via exponentially weighted moving averages and fine-tune model weights through online gradient descent---without full retraining---would enable perpetual operation with quarterly rather than annual recalibration. Federated learning approaches, where multiple buildings collaboratively update a shared model while preserving local privacy, could accelerate adaptation across building portfolios.

From a theoretical perspective, three open questions merit investigation. First, can the variance-entropy criterion be rigorously derived from information-theoretic principles (rate-distortion theory, information bottleneck) rather than empirically optimised via grid search? Such a derivation might reveal optimal weighting functions beyond linear combination and suggest adaptive $\alpha(t)$ schedules that track signal characteristics. Second, how does quantisation error propagate through cascaded sensor-inference pipelines, and can this be formalised through error analysis to guide architectural design? Understanding the interaction between ADC quantisation noise and PTQ rounding error could inform joint optimisation algorithms that minimise end-to-end error rather than independently minimising each stage. Third, what are the fundamental limits of energy-accuracy trade-offs in time-series IoT sensing? Establishing Shannon-theoretic bounds on the minimum energy required to achieve a target classification accuracy under bandwidth and latency constraints would contextualise VABQ's performance relative to theoretical optimality.

\subsection{Concluding Remarks}
\label{sec:conc_remarks}

The proliferation of low-cost IoT sensor networks for environmental monitoring presents a dual imperative: maximising data fidelity to support public health interventions while minimising energy consumption to ensure operational sustainability. This work demonstrates that these objectives need not conflict when temporal signal structure is systematically exploited across both sensing and inference layers. By unifying adaptive ADC resolution control with neural network quantisation under a principled variance-entropy criterion, VABQ achieves substantial energy reductions (51.6\% system-level in duty-cycled deployments) while maintaining classification accuracy above 96\% on real-world indoor air quality data spanning diverse deployment contexts.

The broader implication is that joint optimisation across traditionally separate system layers---sensors, communication, computation---offers a fertile research direction for sustainable IoT. Just as co-design of algorithms and hardware has driven efficiency gains in machine learning accelerators, co-design of sensing and inference pipelines can unlock compound energy savings in resource-constrained edge deployments. As environmental monitoring networks scale from tens to thousands of nodes, and from days to years of continuous operation, such synergistic optimisation will become essential to balance societal benefit with environmental cost.

The VABQ framework, validated on 89.1 million real-world sensor samples, represents a concrete step toward energy-efficient, privacy-preserving, and scalable environmental intelligence. By demonstrating practical feasibility on commodity hardware (ESP32) with manageable calibration overhead (one week per site), this work lowers the barrier to large-scale deployment. The release of the DALTON dataset at NeurIPS 2024 and the planned open-sourcing of the VABQ implementation will enable the research community to build upon these findings, accelerating progress toward sustainable IoT infrastructure for global environmental health monitoring.

%% ============================================================
%% END OF SECTION 7
%% ============================================================
